\paragraph{Gauge parameter and gauge function in the comparison of \(G_T\) with \(G_A\) (conditional; not closure).}
Keep the setting, the class and the normalisation of the two preceding paragraphs. Write \(e_i\) for the contribution of a sector of weight \(w_i\) to the bracket \(1-N_{\mathrm{tot}}/24\) of the preceding paragraph, so that \(\sum_iw_ie_i\) is the net eigenvalue on the two modes with \(\delta\rho=0\) and \(c_T=-4\sum_iw_ie_i\); a scalar component has \(e=-\tfrac1{24}\) there and \(-\tfrac18\) on the mode with \(\delta\rho\neq0\). Let \(D_\mu=\{x_\mu,\cdot\}\), anti-self-adjoint on \(L^2(\rho\,\mathrm d^4x)\), \((\mathrm{Gr}\varphi)_\nu=D_\nu\varphi\), and \(\mathfrak J_\nu=D^\mu\Theta_{\mu\nu}\), the Poisson limit of the tree-level equation of motion; for the perturbation of the quadratic-form paragraph \(\mathfrak J_\nu=-\epsilon\cos(p\cdot x)\,(Mu)_\nu\), with \(M=q^2\eta-qq\) the kernel of Remark~\ref{rem:ch3-trace-quadratic}. The background gauge \(\tfrac1{2\alpha}\rho\,(D_\mu a^\mu)^2\), \(\alpha>0\), \(\kappa=1-1/\alpha\), leaves the ghost operator unchanged and replaces \(\mathcal O_v\) by the non-minimal operator \(\mathcal O_\kappa=\mathcal O_v-\kappa\,\mathrm{Gr}\mathrm{Gr}^\dagger\), with principal symbol \(q^2\mathbb I-\kappa\,qq^{T}\); one proper-time cut-off on \(\mathcal O_\kappa\) serves the effective action and the replica, and the normalisation of the gauge-fixing weight is local and carries no two-derivative term.

For every Poisson bivector \(\mathrm{Gr}^\dagger\mathrm{Gr}=\mathcal O_s\) and \(\mathcal O_v\mathrm{Gr}=\mathrm{Gr}\mathcal O_s-\{\mathfrak J_\nu,\cdot\}\): the commutator terms are fixed by \([D_\mu,D_\nu]=\{\Theta_{\mu\nu},\cdot\}\) and the Leibniz property of \(D_\mu\) on the bracket, and the \(\{\Theta_{\mu\nu},\cdot\}\) coupling of \(\mathcal O_v\) cancels them up to \(\mathfrak J\) (checked in exact arithmetic for a nonlinear bivector). Hence for \(\mathfrak J=0\), \(\mathcal O_\kappa\mathrm{Gr}=(1-\kappa)\mathrm{Gr}\mathcal O_s\), the longitudinal sector is the scalar operator with proper time rescaled by \(1-\kappa\), and up to zero modes \(\operatorname{Tr}e^{-s\mathcal O_\kappa}=\operatorname{Tr}e^{-s\mathcal O_v}+\operatorname{Tr}e^{-s(1-\kappa)\mathcal O_s}-\operatorname{Tr}e^{-s\mathcal O_s}\). The integrated \(\LUV^2\) density of the connection sector is then \(\tfrac{3+\alpha}6\,\rho R[g_\rho]+\rho^2X\) on such backgrounds, with the coefficient \(1\) of \(\rho^2X\) independent of \(\alpha\). On the replica cover with constant \(\bar\Theta\) the same intertwining, with the domain at \(\Sigma\) for which the longitudinal sector is the image of the regular scalar domain, gives the area coefficient \(\tfrac{3+\alpha}6\) for the connection sector in the units \(c\) of the table above, equal to the coefficient of \(\rho R[g_\rho]\). This depends on that choice of domain; other choices are not examined.

About constant \(\bar\Theta\), with \(\mathfrak J\neq0\), a numerical second-order evaluation of the heat trace (momentum-space Duhamel expansion, Gauss--Hermite quadrature, contour extraction of the \(p^4\) coefficient; it reproduces the scalar value above, the \(\alpha=1\) values \(\tfrac56\) and \(\tfrac12\) of the connection sector, and the weight \(3+\alpha^2\) of the \(\LUV^4\) term) gives for the connection sector the values, matched by the rational forms, \(\tfrac56-\tfrac{\alpha-1}{24}\) on the two modes with \(\delta\rho=0\) and \(\tfrac12-\tfrac{\alpha-1}8+\tfrac{\alpha^2-1}4\) on the mode with \(\delta\rho\neq0\) (\(\alpha=\tfrac13,\tfrac12,\tfrac23,1,2,3\) in one background, \(\tfrac12,1,2\) in another; rational fits to \(2\times10^{-3}\), not derived; \(\alpha=\tfrac13,\dots,3\) only). The two modes with \(\delta\rho=0\) thus show the \(\alpha\)-dependence of the \(\mathfrak J=0\) result, and the remainder, \(\tfrac{\alpha^2-1}2\rho|\nabla\ln\rho|^2_{g_\rho}\), acts on the mode with \(\delta\rho\neq0\) only. With scalar and ghost sectors independent of \(\alpha\), \(c_T(\alpha)=\tfrac{N_s+\alpha-1}6-\tfrac{11}3\) and the replica coefficient is \(c_A(\alpha)=\tfrac{N_s+\alpha+1}6\), so \(1/G_A-1/G_T=2\LUV^2/\pi\) for the tested range \(\tfrac13\le\alpha\le3\); on backgrounds with \(\mathfrak J=0\) and under the conditions above the \(\alpha\)-dependence cancels in the difference for every \(\alpha>0\). The TT value agrees with the exact lemma on \(\mathfrak J=0\) backgrounds. For \(\alpha\neq1\) the form is not spanned by \(\rho R[g_\rho]\) and \(\rho^2X\); the third-mode to \(\delta\rho=0\) ratio of a term \(c\,\rho R[g_\rho]\) is \(3\), and the net form attains this ratio at \(\alpha=3\), where the gap is unchanged, so that ratio does not diagnose the coefficient of \(\rho^2X\) outside the Feynman gauge. A gauge function \(D_\mu a^\mu+c\,\mathfrak J_\mu a^\mu\) (\(\alpha=1\), ghost operator \(-D_\mu D_\mu-c\mathfrak J_\mu D_\mu\)) adds \(c^2\rho^2\mathfrak J_\nu\mathfrak J_\nu\) to the density of the connection and ghost sectors together at quadratic order about \(\bar\Theta\), that is \(c^2q^2/(2p_\eta^2)\) to the eigenvalue on all three modes (two backgrounds, \(c=\tfrac12,1,2\)); this lowers \(c_T\) at fixed \(c_A\) and depends on the direction of \(p\) unless \(\bar\Theta\) is isotropic. Gauge functions with other first-derivative vectors are not examined.

\paragraph{Fermionic content and the connection-sector cone in the comparison of \(G_T\) with \(G_A\) (conditional; not closure).}
The displayed action is bosonic. Add \(n_D\) Dirac fields with action \(\int\rho\,\bar\Psi\gamma^\mu D_\mu\Psi\), the Poisson limit of \(\operatorname{Tr}\bar\Psi\Gamma^\mu[\hat X_\mu,\Psi]\), and the connection gauge \(\alpha=1\); the Yukawa couplings \(\Gamma^i[\hat X_i,\Psi]\) to the \(N_s\) transverse matrices vanish on the background and are dropped. The operator \(\mathcal O_f=-(\gamma^\mu D_\mu)^2=\mathcal O_s-\tfrac12\gamma^{\mu\nu}\{\Theta_{\mu\nu},\cdot\}\) has the form~\eqref{eq:O-eta} on \(G\) with \(\omega_\sigma=-\tfrac12\partial_\sigma\ln\rho\,\mathbb I+\beta_\sigma\), \(\beta^\sigma=\tfrac14\gamma^{\mu\nu}\Theta^{\rho\sigma}\partial_\rho\Theta_{\mu\nu}\), the spinor lift of \(\mathcal B^\sigma\), and \(\operatorname{tr}\beta_\sigma\beta^\sigma=-\tfrac12X\) over four components (the vector value is \(-X\)); it contains no explicit curvature, so \(\sqrt{|G|}\operatorname{tr}(\tfrac16R_G-E_f)=\tfrac46\rho R[g_\rho]+\tfrac12\rho^2X\), with eigenvalues \(\tfrac13,\tfrac13,0\) on the three modes (two backgrounds). The statistical weight is \(-1\) per Dirac field (\(\Gamma=-\tfrac12\ln\det\mathcal O_f\) over four complex components). The coefficient of \(\rho^2X\) of the whole content is therefore \(1-n_D/2\): the weight of a field is its index under the rotation generated by \(\mathcal B_\sigma\), and it vanishes for \(n_D=2\), the sixteen real components of the ten-dimensional maximally supersymmetric matrix model~\cite{IKKT1997} with the Poisson background along four coordinates and \(N_s=6\). No symmetry of the displayed action forces this. For \(n_D\) fields \(c_T=\tfrac{N_s+\alpha-1}6+\tfrac{4n_D}3-\tfrac{11}3\). With the antiperiodic polar-frame condition, which gives the Cartesian-frame components the phase \((-1)^{n+1}\), the area coefficient of a Dirac field is that of two scalar components, \(+\tfrac13\) (Equations~(126)--(127) of~\cite{Solodukhin2011}; the zeta-function evaluation of the twisted cone gives \(-\tfrac12\) times the scalar value per component on the branch \(n<1\)). Then \(c_A=\tfrac{N_s+\alpha+1}6+\tfrac{n_D}3\) and \(c_A-c_T=4-n_D\). This gluing transports the spinor index with the replica rotation; it is not the unrotated pull-back gluing used above for \(a^\mu\). With unrotated periodic gluing of the Cartesian components the area coefficient of a Dirac field is \(-\tfrac23\), equal to its coefficient of \(\rho R[g_\rho]\), and \(c_A-c_T=4-2n_D\), which vanishes at \(n_D=2\) with \(c_T=c_A=0\): an equality of vanishing \(\LUV^2\) coefficients, not a derivation. The \(\rho^2X\) coefficient vanishes at \(n_D=2\); for antiperiodic gluing \(c_A-c_T=2\) there and equality needs \(n_D=4\), where \(\rho^2X\) has coefficient \(-1\). The gluing selected by the matrix state is not determined here. For \(n_D=2\), \(N_s=6\) one has \(N_s+2-4n_D=0\), \(c_T=0\) and \(c_A=2\): the \(\LUV^4\) and \(\LUV^2\) tensor coefficients vanish, and the area coefficient is \(2\) for antiperiodic and \(0\) for periodic gluing.

If the internal index of \(a^\mu\) is rotated with the replica angle, as for a vacuum invariant under the combined rotation, the area coefficient of the connection sector at \(\alpha=1\) is \(-\tfrac13\) on the branch \(n<1\) and \(\tfrac53\) on \(n>1\) instead of \(\tfrac23\); with the ghosts these are \(-\tfrac23\), which is \((D_1(4)-6)/6\) of Equation~(147) of~\cite{Solodukhin2011}, and \(\tfrac43\), instead of \(\tfrac13\); this function of \(n\) is not analytic at \(n=1\), and the gap is \(3-n_D\) or \(5-n_D\). These statements are for the Poisson-limit operators, the \(\LUV^2\) term, one loop, the quadratic form about constant \(\bar\Theta\) on the two modes with \(\delta\rho=0\), the gauge families and the replica prescriptions stated. They do not derive the content, the fermion coupling, the gluing or the rotation of the internal index, which belong to Problem~A, Assumption~\ref{ass:EE} or the modular flow; they do not treat Moyal corrections, gauge functions beyond those stated, \(\Phi\)--\(a\) mixing, or the Kaluza--Klein sector. Problems~B, C, D, F and~G remain open.
